\section{Scales}

A trit is not a field, and a dock is not a particle. Recognizable physics is the coarse-grained face the tones show at the cusp, the smooth average of a great many of them rather than any single one, and to read one tone as a field or a particle is the commonest way to fool yourself. The principle of two layers, stated at the start, becomes a working method here.

There is no choice about this, only the question of doing it well. A single self is a state drawn from three raised to twenty-four billion, and the universe is past all such numbers, so reading physics off the base dock by dock and beat by beat is not merely hard but impossible, now and for any machine there will ever be. The only way to do physics at all is to read at the right rung, the fine law where it matters and the blurred effective law everywhere else, the way no one computes a magnet from every atomic spin but reads its smooth field, or a gas from every molecule but reads its temperature and pressure. Levels of detail are not a convenience here, they are the one available method, the same one every science already runs on.

Coarse-graining means averaging, and a magnet shows it plainly. Each atom of a magnet carries a spin that flips in whole jumps, up or down, nothing between. Yet the magnet has a smooth field that points any direction and varies by any amount. The smoothness is not in the atoms. It is in the average of many of them across a patch. The same move turns the base into physics. Average the ternary tones over a region and the average is a fine-grained quantity, effectively a real number, a field. Momentum, charge, and current all appear this way, as averaged tones, never as single ones. This is also what every decimal in the results is. The anisotropy of $0.059$, the Bell value of $2.83$, the gyromagnetic ratio of two, none is a number in the law. Each is a measurement of the integer dynamics, the way a temperature is read off a gas of whole atoms. The reals are in the ruler, never in the thing measured. The smooth world is a song the trits sing, and only the singers are real.

Doing this step after step is renormalization, the climb from fine to coarse, and at each rung some detail washes out while a few stubborn quantities survive. The survivors are the physics. The base has no electron and no photon inside it. What it has is a law whose coarse-grained survivors, a rung or two up, move like an electron and a photon. The whole of the physics part is a catalog of survivors.

The climb does more than simplify, it makes new things. Averaging electrons and nuclei gives the atom, averaging atoms gives the unit of sound, averaging quarks gives the proton, averaging neurons gives the thought, and each is a real object at its own rung, often steadier and more telling than the swarm beneath it. So the question the base poses is the one modern physics has been turning toward, not what the smallest pieces are but what the right emergent variables are, the slow survivors, the conserved quantities, the stable patterns that carry the next level up. And the standard the program holds itself to is exactly the one that certifies emergence everywhere else. Run the microscopic rule forward, coarse-grain, and check that the higher law appears on its own, never slipped in by hand. Where it appears the result counts, and where it does not the negative is kept.

That the climb is consistent was checked sharply. Coarse-grain the same configuration at blocks of two, four, eight, and sixteen, and the charge stays exact at every level while the effective coupling flows under the standard block-spin recursion to a fixed point, the value it should reach. A rule learned at the coarse scale reproduces the fine dynamics when the system is ordered and fails, as it must, when it is frustrated, the clean sign that the macro law is real rather than imposed. That physics is the same kind of thing at every scale, the deep reason coarse-graining works at all, is here a measured fact and not an assumption.

Not every climb succeeds, and the failures are reported. One natural way to coarse-grain the bulk, averaging in shells outward from a point, does not produce the layered tower of scales that real physics shows. We built it, measured it, and found no tower. \cite{pollard2026vibetest} That is a negative, kept in plain view, and it marks the coarse-graining of the full bulk as an open problem rather than a solved one. The flat cusp coarse-grains cleanly. The curved bulk does not yet.

The lesson governs every claim that follows. A result counts only when it is read at the right rung. A clustering of raw tones is not gravity, and a single high tone is not a particle, however suggestive it looks. The physics is what survives the climb, and the climb is part of the measurement.
